\documentclass[11pt,a4paper]{report}

\usepackage[margin=2.6cm]{geometry}
\usepackage{listings}
\usepackage{xcolor}
\usepackage{booktabs}
\usepackage{microtype}
\usepackage{amsmath}
\usepackage{amssymb}
\usepackage{tikz}
\usepackage{longtable}
\usepackage[colorlinks=true,linkcolor=black,urlcolor=blue,citecolor=black]{hyperref}
\usetikzlibrary{positioning,arrows.meta}

% long \code{} tokens cannot hyphenate; let TeX stretch rather than overflow
\setlength{\emergencystretch}{2.5em}

\definecolor{kw}{RGB}{0,80,160}
\definecolor{cmt}{RGB}{90,110,90}
\definecolor{str}{RGB}{150,60,40}
\definecolor{rule}{RGB}{200,200,200}
\definecolor{shade}{RGB}{248,248,246}

\lstset{
  language=C++,
  basicstyle=\ttfamily\small,
  keywordstyle=\color{kw}\bfseries,
  commentstyle=\color{cmt}\itshape,
  stringstyle=\color{str},
  numbers=none,
  frame=single,
  rulecolor=\color{rule},
  backgroundcolor=\color{shade},
  breaklines=true,
  breakatwhitespace=true,
  showstringspaces=false,
  columns=fullflexible,
  keepspaces=true,
  captionpos=b,
  xleftmargin=0pt,
  framexleftmargin=3pt,
  morekeywords={active,largeint,Tape,Jacobian,TapeState,Process,Vertex,Edge,
                AdjEntry,AdjArena,Adjacency,EdgeArena,BreakException,elim_t,
                thread_local,noexcept,override,nullptr}
}

\newcommand{\svegp}[1]{\textsc{svegp-#1}}
\newcommand{\code}[1]{\texttt{#1}}

\title{\vspace{-2cm}\Huge Maxwell\\[0.3cm]
\Large Automatic Differentiation by Sparse Graph Vertex Elimination\\[0.6cm]
\large Implementation Report and User Guide}
\author{Software and Tools for Scientific Computing}
\date{9 September 2026}

\begin{document}

\maketitle

\begin{abstract}
\noindent
Maxwell is a C\texttt{++} operator-overloading library for automatic
differentiation. It records a program's arithmetic as a directed acyclic graph
and computes the Jacobian by \emph{vertex elimination} rather than by a forward
or reverse sweep, which lets the elimination direction be chosen per problem
and makes the cost of differentiation a property of the graph rather than of
the traversal. Its distinguishing feature is \emph{checkpointing under a memory
budget}: a tape too large to hold is broken into partitions, and the recorded
section is replayed once per partition, so a deep computation can be
differentiated in bounded memory. The partitioning is bit-exact --- the same
Jacobian is produced at a 1\,GB budget and at a 1\,kB budget, to the last bit,
with twelve pinned witnesses asserting it.

This report documents the implementation and the reasoning behind it. Part~I is
a maintainer's reference: the graph representation, the elimination core, the
checkpoint protocol, the memory accounting, the tape object, the RAII surface,
the two storage arenas, and what is and is not possible concurrently. Design
choices are justified with the measurements that produced them, including
several that were measured \emph{wrong} before they were measured right ---
a memory model that under-reported by a factor of $2.44$, an arena whose first
two parameter choices each made things worse, and a scope-restoration bug that
only a constructed test could expose. Part~II is a user guide: how to drive the
library, what the checkpoint contract requires of your code, how to choose a
budget, and six worked use cases taken from the shipped examples.

\vspace{0.4cm}
\noindent
All measurements were taken on GCC 13.3.0, \code{-O2}, x86-64, and are
reproducible from the tree with \code{make test} and \code{make sanitize}.
\end{abstract}

\tableofcontents

\part{Design and Implementation}

\chapter{Introduction}

\section{What the library does}

Maxwell differentiates a C\texttt{++} program by overloading the arithmetic
operators on a scalar type, \code{maxwell::active}. Where a program computes
with \code{double}, it computes with \code{active} instead; each elementary
operation appends a vertex and one or two edges to a graph whose edge weights
are the local partial derivatives. Once the program has run, the Jacobian of
the declared dependents with respect to the declared independents is obtained
by eliminating every intermediate vertex from that graph.

The public protocol, in its oldest form, is seven ordered steps:

\begin{lstlisting}
initialize(n,m,budget);
  independent(x[i]) ...
  while( checkpoint(x,y) ){ try{ f(); } catch(BreakException const&){} }
  dependent(y[k]) ...
  harvest(m,n,A);          // double**, caller-owned
  free_jacobian(m,A);
finalize();
\end{lstlisting}

\noindent
Part~II describes the modern surface, in which a \code{Tape} object replaces
\code{initialize}/\code{finalize}, owns the loop and its \code{try}/\code{catch},
and returns the Jacobian as a value. The free functions above remain, are still
public, and are what the object surface is built on --- Chapter~\ref{ch:raii}
explains why they were kept rather than retired.

\section{What makes it different}

Two things.

\paragraph{Vertex elimination, not a sweep.} Forward-mode and reverse-mode AD
are the two extreme orderings of a more general operation: eliminating vertices
from the linearised computational graph. Maxwell implements the general
operation, so the direction is a setting (\code{set\_elim\_mode}) rather than a
different code path. On a deep, narrow tape the difference is large and
measurable: \code{examples/heat\_assim} reports $6\,835$ multiply--adds in
reverse against $112\,102$ forward on the same graph, a ratio of $16.4\times$,
for an identical Jacobian.

\paragraph{Checkpointing under a byte budget.} This is the feature the rest of
the design serves. The library is told how many bytes of graph it may hold. It
runs the user's section once to \emph{profile} it, discovers how many partitions
the budget implies, and then replays the section once per partition, recording
only one partition at a time and eliminating each into the accumulating result.
Deep time-stepping problems that would otherwise exhaust memory become
tractable, and --- crucially --- the answer does not depend on the budget. The
regression suite asserts this at seven budgets from 200\,000 down to 1\,000
bytes with \emph{zero} deviation, not a tolerance.

\section{A note on the listings}

Code in this report is taken from the tree, not paraphrased. Listings are
excerpted --- elisions are marked \code{...} --- and occasionally carry a
short explanatory comment that is not in the source. Where a listing shows a
defect it is labelled as such and the fixed form is given alongside. Every
snippet was checked against the file it came from before this document was
compiled.

\section{Scope of this report}

Chapter~\ref{ch:graph} covers the recording machinery and the graph
representation. Chapter~\ref{ch:elim} covers elimination. Chapter~\ref{ch:ckpt}
covers checkpointing and the contract it imposes on user code.
Chapter~\ref{ch:mem} covers the memory accounting that the budget is spent
against. Chapter~\ref{ch:tape} covers the tape object that replaced a
singleton. Chapter~\ref{ch:raii} covers the RAII surface.
Chapter~\ref{ch:arena} covers the two storage arenas and the measurements that
shaped them. Chapter~\ref{ch:threads} covers concurrency.
Chapter~\ref{ch:verify} covers how the library is verified, which is a design
choice in its own right.

\section{The source tree}

\begin{center}
\begin{tabular}{@{}ll@{}}
\toprule
\code{maxwell.hpp}            & the public surface; includes nothing internal \\
\code{inc/Active.hpp}         & the \code{active} scalar and its operators \\
\code{inc/API.hpp}            & the free-function protocol \\
\code{inc/Tape.hpp}           & \code{maxwell::Tape} and \code{maxwell::Jacobian} (public) \\
\code{inc/TapeState.hpp}      & \code{internals::TapeState} --- one tape's whole state \\
\code{inc/Process.hpp}        & the recorder and eliminator \\
\code{inc/Vertex.hpp}         & \code{Vertex}, \code{Adjacency}, \code{AdjArena} \\
\code{inc/Edge.hpp}           & \code{Edge}, \code{EdgeArena} \\
\code{inc/Typedefs.hpp}       & \code{largeint}, \code{elim\_t} \\
\code{src/*.cpp}              & six translation units \\
\code{examples/}              & nine programs; \code{make test} runs all of them \\
\bottomrule
\end{tabular}
\end{center}

\noindent
\code{maxwell.hpp} deliberately does not include \code{Vertex.hpp},
\code{Edge.hpp}, \code{Process.hpp} or \code{TapeState.hpp}. That is not
tidiness: it is what allowed the graph representation to be replaced twice
(Chapter~\ref{ch:arena}) without a single caller changing, and it is why the
public \code{Tape} holds one forward-declared pointer rather than its state by
value.

\chapter{Recording: the graph}
\label{ch:graph}

\section{The active scalar}

\code{active} is a 48-byte value carrying its numeric value and its position in
the graph:

\begin{lstlisting}
class active
{
public:
  mutable bool     reachable;   // does this value depend on an independent?
  mutable largeint idx;         // vertex index, 0 if passive
  mutable largeint owner_idx;   // which partition recorded it
  mutable largeint old_idx;     // saved across a replay
  mutable double   val;         // the value itself
  mutable internals::Vertex * vtx;   // cached vertex, or NULL
  ...
};
\end{lstlisting}

Every member is \code{mutable} so that the operators can take
\code{const active\&} and still update bookkeeping. This is defensible for
\code{idx} and \code{vtx}, which are cache-like state on a logically constant
object; it is less defensible for \code{val} and \code{reachable}, which are
the object's actual value, and it means \code{const active\&} promises less
than it appears to. It is recorded here as a known wart rather than defended.

\paragraph{\code{active} is not a value type.} Copying one \emph{records}. The
copy constructor is:

\begin{lstlisting}
active::active( const active & x ):
reachable(x.reachable), idx(0), owner_idx(0), old_idx(0), val(x.val), vtx(NULL)
{
  unary_op( x , 1.0 , *this , false );   // an identity vertex and edge
}
\end{lstlisting}

ADOL-C and CoDiPack do the same, and it is the only way an overloading library
can track a copy. But nothing in the type says so, and the standard library
copies things constantly. Measured on \code{y = x0*x1 + x0}:

\begin{center}
\begin{tabular}{@{}lr@{}}
\toprule
what the program does & tape size \\
\midrule
the straight-line expression                          & 480 bytes \\
\dots with one \code{active} passed \emph{by value}   & 600 bytes \\
\dots with 8 \code{push\_back} into \code{std::vector<active>} & \textbf{2\,640 bytes} \\
\bottomrule
\end{tabular}
\end{center}

\noindent
Eight \code{push\_back}s into a growing vector cause $1+2+4+8=15$ reallocation
copies, and every copy records an identity vertex and edge: exactly the
$15\times144$ bytes measured. The derivative is still correct; the tape is five
and a half times larger than it should be, and nothing warns. Section~\ref{sec:pitfalls}
tells users how to avoid it; Section~\ref{sec:independents} describes the
library-side mitigation, which is that the tape now hands out its own
independents from storage sized exactly once.

\section{How an operation is recorded}

Each overloaded operator computes the value and the local partials, then calls
one of five recording entry points. Multiplication is representative:

\begin{lstlisting}
active maxwell::operator*( const active & x1 , const active & x2 )
{
  active x3;
  x3.val = x1.val*x2.val;
  x3.reachable = x1.reachable || x2.reachable;
  binary_op( x1 , x2.val , x2 , x1.val , x3 );   // dy/dx1 = x2, dy/dx2 = x1
  return x3;
}
\end{lstlisting}

\code{binary\_op} reaches the current tape and forwards to \code{Process}, whose
behaviour differs sharply between the two passes:

\begin{lstlisting}
void Process::binary_op( const active & x1 , double dy_dx1 ,
                         const active & x2 , double dy_dx2 , const active & x3 )
{
  if(profiling){
    if(x3.reachable){
      x3.idx = next_vertex_idx++;
      intmed_count++;
      if(&x1==&x2){ edge_count++; }
      else{
        if(x1.reachable) edge_count++;
        if(x2.reachable) edge_count++;
      }
    }
  }else{
    Vertex * v1=NULL , * v2=NULL , * v3=NULL;   // SVEGP-20
    ...
    if(x3.reachable){
      x3.idx = next_vertex_idx++;
      v3 = vertex_on_lhs(x3);
      if(&x1==&x2){ add_edge(v1,v3,dy_dx1+dy_dx2); edge_count++; }
      else{
        if(x1.reachable){ add_edge(v1,v3,dy_dx1); edge_count++; }
        if(x2.reachable){ add_edge(v2,v3,dy_dx2); edge_count++; }
      }
    }
  }
  check_memory();
}
\end{lstlisting}

Three things in that fragment are load-bearing.

\paragraph{The profiling pass allocates nothing.} It only counts. That is what
makes the budget decision possible before a single vertex exists, and it is
what Chapter~\ref{ch:arena} exploits: the exact vertex and edge counts are known
in advance.

\paragraph{Aliasing is handled explicitly.} \code{\&x1==\&x2} detects
\code{x*x}, where a naive implementation would add two edges between the same
pair of vertices and then have to merge them. The partials are summed instead.
\code{examples/invariants} pins this path as its \code{alias} case.

\paragraph{The \code{NULL} initialisers are a fix, not style.} \svegp{20}:
those pointers were left uninitialised, and \code{add\_edge} only guards against
\code{NULL}. An indeterminate non-null value walked straight into
\code{tgt->add\_in\_edge(src,\dots)}. The library's own callers always set the
reachability flags consistently, so it never fired internally --- but
\code{binary\_op} and its siblings are \emph{exported} for hand-written
operators, and there a mismatched flag turns a dropped edge into a wild pointer
dereference.

\section{Reachability}

\code{reachable} is the flag that keeps passive arithmetic out of the graph. A
value is reachable if it descends from a registered independent. A computation
on two passive values records nothing: \code{operator=} from a \code{double}
calls \code{passive\_op}, which kills the vertex and sets \code{idx=0}. This is
why a program can mix \code{active} and \code{double} freely and only pay for
the part that actually depends on the inputs.

\section{The graph}

A \code{Vertex} is 72 bytes and holds nothing but its identity and its two
adjacency lists:

\begin{lstlisting}
class Vertex
{
public:
  bool     alive;       // true if this vertex is a declared dependent
  largeint idx;
  largeint owner_idx;

  Adjacency in_edges;
  Adjacency out_edges;
  ...
};
\end{lstlisting}

An \code{Edge} is 24 bytes --- a partial derivative and its two endpoints:

\begin{lstlisting}
class Edge
{
public:
  double   eval;    // the local partial derivative
  Vertex * src;
  Vertex * tgt;
  ...
};
\end{lstlisting}

\code{alive} deserves a note, because the name is misleading: it does not mean
``not yet destroyed''. It marks a vertex as a \emph{declared dependent}, and its
only effect is that elimination refuses to eliminate it. Dependents must survive
because they are the rows of the Jacobian.

Adjacency is a sorted array of \code{(neighbour index, edge pointer)} pairs,
binary searched. Its members are called \code{first} and \code{second} because
the elimination code was written against \code{std::map} and is full of
\code{it->second}; keeping the names meant the representation could be replaced
without touching a line of the algorithm. Chapter~\ref{ch:arena} covers how the
storage beneath it evolved.

\chapter{Elimination}
\label{ch:elim}

\section{The operation}

Eliminating an intermediate vertex $v$ removes it from the graph while
preserving every path through it. For each in-edge $(u,v)$ with weight $a$ and
each out-edge $(v,w)$ with weight $b$, the product $ab$ is added to the edge
$(u,w)$, creating it if it does not exist. The cost is the number of such
products:
\[
  \mathrm{cost}(v) = \deg^{-}(v)\cdot\deg^{+}(v).
\]

The implementation is the whole of \code{Vertex::eliminate}:

\begin{lstlisting}
largeint Vertex::eliminate( EdgeArena & arena , AdjArena & adj )
{
  largeint m = in_edges.size();
  largeint n = out_edges.size();
  largeint cost = m*n;

  for( outedge_it=out_edges.begin() ; outedge_it!=out_edges.end() ; outedge_it++ )
    outedge_it->second->tgt->in_edges.erase(this->idx);

  for( inedge_it=in_edges.begin() ; inedge_it!=in_edges.end() ; inedge_it++ )
  {
    inedge_it->second->src->out_edges.erase(this->idx);

    for( outedge_it=out_edges.begin() ; outedge_it!=out_edges.end() ; outedge_it++ )
    {
      double cij = (inedge_it->second->eval)*(outedge_it->second->eval);

      Edge * direct_link = NULL;
      Adjacency::iterator direct_link_it =
        inedge_it->second->src->out_edges.find(outedge_it->second->tgt->idx);

      if(direct_link_it!=inedge_it->second->src->out_edges.end())
        direct_link = direct_link_it->second;

      if(direct_link){
        direct_link->eval += cij;                       // fold into an existing edge
      }else{
        outedge_it->second->tgt->add_in_edge(            // fill-in
          inedge_it->second->src , cij , arena , adj );
      }
    }
    arena.release( inedge_it->second );
  }

  for( outedge_it=out_edges.begin() ; outedge_it!=out_edges.end() ; outedge_it++ )
    arena.release( outedge_it->second );

  in_edges.clear(adj);
  out_edges.clear(adj);
  return cost;
}
\end{lstlisting}

The \code{find} on the predecessor's out-list is the \emph{direct-link probe},
and it is the reason adjacency must stay sorted rather than becoming a linked
list. Measured degrees on that list are 40 entries at $n_x=20$ and 109 at
$n_x=40$, growing with problem size; seven comparisons over contiguous memory
beat a tree walk, but a linear scan over 109 scattered nodes would have made
elimination asymptotically worse on exactly the deep tapes the library exists
for.

\section{Order matters, and adjacency order does not}

This distinction was assumed backwards during design and was settled by
experiment.

Within one \code{eliminate()} call, \code{in\_edges} is keyed by source index
and \code{out\_edges} by target index, so each $(u,w)$ pair occurs exactly once
and \code{direct\_link->eval += cij} fires at most once per pair per
elimination. All accumulation \emph{across} eliminations is ordered by
elimination order, which comes from \code{intmed\_vec} and \code{intmed\_map},
never from adjacency.

Reversing the adjacency iteration order therefore changes nothing --- verified
by doing it, against twelve pinned Jacobian hashes, all of which stayed
byte-identical. That freed the storage layout to be chosen for speed alone.

Elimination order, by contrast, changes the floating-point result in the last
bits, exactly as it should: \code{examples/invariants} records
\code{nonlinear/rev} at hash \code{0xaeb4c2dad9593bc3} and \code{nonlinear/fwd}
at \code{0x814891dd92fad4e7}, with costs 3\,291 and 3\,333. Same Jacobian to
within rounding; different accumulation order; different last bits. Any change
that reorders elimination is therefore a change to the numbers, and the pinned
table will say so.

\section{Direction}

\begin{lstlisting}
largeint Process::eliminate()
{
  if(elim_mode==FORWARD_ELIM){ return forward_eliminate(); }
  return reverse_eliminate();
}
\end{lstlisting}

Forward elimination walks \code{intmed\_vec} in order; reverse walks it
backwards. Both refuse to eliminate a vertex whose \code{alive} flag is set, and
both skip vertices with no in-edges (independents) or no out-edges that are
dependents.

\svegp{23} is worth recording: the direction was not selectable at all.
\code{checkpoint()} called \code{reverse\_eliminate()} unconditionally, so
forward elimination could never be the elimination that produces the Jacobian
--- calling \code{forward\_elimination()} afterwards merely walked an
already-eliminated graph and reported the same cost. The \code{elim\_t} enum had
sat unused in \code{Typedefs.hpp} since the beginning, and the old test Makefile
recorded that a \code{set\_elim\_mode()} had once been \emph{removed} from the
API, so restoring it was recovering an intended feature rather than inventing
one.

\section{Harvest}

After elimination, the surviving edges from independents to dependents are the
Jacobian entries. \code{Process::harvest} walks the dependents, and for each
looks up the edge from each independent:

\begin{lstlisting}
Edge * e = dep_vPtr->from(indep_vPtr);

if(e){
  A[dep_pos][i] = e->get_partial();
}
else if(indep_vPtr==dep_vPtr){
  A[dep_pos][i] = 1.0;     // a dependent that IS an independent
}
\end{lstlisting}

The \code{indep\_vPtr==dep\_vPtr} branch handles \code{y = x} --- a dependent
that is also an independent has no edge to itself and a derivative of exactly
one. \code{examples/invariants} pins this as its \code{passthrough} case.

Two defects lived here. \svegp{03}: the column of an entry is the
\emph{position} of the independent, not its global vertex index; using
\code{get\_idx()-1} coincides with the column only while the independents happen
to occupy indices $1..n$ in registration order, and registering one after any
active arithmetic both lost a column and wrote past the end of the row.
\svegp{05}: the requested shape used to be trusted, so a caller passing $m$ and
$n$ transposed overran the heap. It is checked now, and refuses:

\begin{lstlisting}
if( m != this->dep_vec.size() || n != this->indep_count ){
  std::cerr << "maxwell::harvest: requested " << m << "x" << n
            << " but the graph holds " << this->dep_vec.size() << "x" << this->indep_count
            << " -- refusing to write out of bounds.\n";
  return;
}
\end{lstlisting}

\chapter{Checkpointing}
\label{ch:ckpt}

\section{The idea}

A tape for a deep computation does not fit in memory. The classical remedy is
to store the program's intermediate \emph{state} at checkpoints and recompute
between them. Maxwell does the dual: it stores nothing and \emph{replays the
program}, recording a different slice of the graph each time.

The user's computation is placed in a \emph{section} driven by a loop:

\begin{lstlisting}
while( checkpoint(x,y) ){
  try{ f(); }
  catch( BreakException const & ){}
}
\end{lstlisting}

\code{checkpoint} returns \code{true} while another pass is needed. The first
pass is the \emph{profiling} pass: nothing is allocated, everything is counted,
and the partition boundaries are discovered. Then one \emph{productive} pass per
partition follows.

\section{Partitions and the owner index}

Partition membership is carried on two counters in \code{Process}:
\code{next\_owner\_idx}, which advances as the section runs, and
\code{tgt\_owner\_idx}, which names the partition this pass is to record.
Recording is gated by a single scalar equality:

\begin{lstlisting}
bool Process::is_proc()
{
  if(next_owner_idx==tgt_owner_idx){ return true; }
  return false;
}
\end{lstlisting}

\noindent
Every vertex-creating helper begins with \code{if(is\_proc())}. So during a
productive pass the section's arithmetic runs in full, but only the operations
belonging to one partition build graph. \textbf{Exactly one partition is active
at a time, by construction} --- a fact with consequences for parallelism that
Chapter~\ref{ch:threads} returns to.

Boundaries are found by \code{check\_memory}, called after every recorded
operation:

\begin{lstlisting}
void Process::check_memory()
{
  largeint cur_mem_usage = (intmed_count*VERTEX_BYTES + edge_count*EDGE_BYTES);
  largeint part_size = cur_mem_usage - prev_mem_usage;

  if(part_size>=mem_size){
    if(!profiling){
      if(is_proc()){
        tgt_owner_idx--;
        ... move intmed_vec into intmed_map, merging duplicates ...
        if(throwable){ throw BreakException(); }
      }
    }
    prev_mem_usage = cur_mem_usage;
    next_owner_idx++;
  }
}
\end{lstlisting}

When the active partition is complete, the pass has no further use for the rest
of the section, so it is abandoned by throwing \code{BreakException} from inside
the operator. This is the mechanism that makes replay cheap: a pass that has
already recorded its partition does not run the remaining arithmetic at all.

The invariant $\text{BreakExceptions} = \text{chunks} - 1$ holds throughout,
because the first productive pass never throws: \code{reinitialize()} sets
\code{throwable} only once its own run counter is non-zero.

\section{The replay loop}

\begin{lstlisting}
bool maxwell::checkpoint( active * x , active * y )
{
  TapeState * tp = current_tape();
  if(!tp){ /* diagnose and end the loop */ return false; }

  Process * myProc = &tp->proc;

  if(!tp->run_counter){ tp->run_counter++; return true; }   // profiling pass

  if(tp->run_counter==1){
    for( largeint k=0 ; k<tp->dependent_size ; k++ ){
      y[k].old_idx = y[k].idx;
      tp->dep_shadow_copy[k] = y[k].val;
    }
    myProc->finalize();                        // close profiling
    tp->run_target = myProc->get_tgt_owner_idx();
    myProc->disable_profiling();
  }else{
    myProc->eliminate();                       // fold the partition just recorded
  }

  internals::restore_values(x,y);              // put the independents back
  myProc->reinitialize();

  if(tp->run_counter<=tp->run_target){ tp->run_counter++; return true; }

  for( largeint k=0 ; k<tp->dependent_size ; k++ ){
    y[k].idx = y[k].old_idx;
    y[k].val = tp->dep_shadow_copy[k];
  }
  return false;
}
\end{lstlisting}

\code{restore\_values} resets the independents from \code{indep\_shadow\_copy},
which is why the section produces identical values on every pass. Note that
\code{tgt\_owner\_idx} counts \emph{down}: the first productive pass records the
\emph{last} partition, and elimination accumulates backwards through the tape.

\section{The contract, which is the sharpest edge in the library}
\label{sec:contract}

The section is run $1+N$ times, where $N$ is not known until profiling finishes,
and it may be aborted at an arbitrary operator by a \code{throw} from inside
\code{operator*}. Three requirements follow, and for a long time none of them
was written down anywhere:

\begin{enumerate}
\item it must be \textbf{re-runnable} --- same inputs, same result, every time;
\item it must be \textbf{side-effect free} --- the profiling pass runs it an
      extra time, so any I/O or RNG inside it happens once more than the author
      expects;
\item it must be \textbf{exception-safe} --- a \code{BreakException} unwinds
      from the middle of the computation.
\end{enumerate}

Requirement (3) was broken during a port: \code{heat\_assim} originally held its
two state buffers in raw \code{new[]} and swapped them by pointer, leaking
$n_x+1$ actives on \emph{every} break --- thousands across a conjugate-gradient
run. It took a rewrite to \code{std::vector<active>} to fix, and nothing would
have caught it.

The RAII surface (Chapter~\ref{ch:raii}) takes requirement (3) off the caller
entirely by owning the \code{try}/\code{catch}. It cannot take (1) and (2), and
does not pretend to.

\section{The property that makes it trustworthy}

Chunking must not change the answer. \code{examples/regress} asserts this
directly, sweeping budgets against a reference computed at $2^{30}$ bytes:

\begin{center}
\begin{tabular}{@{}lr@{}}
\toprule
budget (bytes) & max relative deviation \\
\midrule
200\,000 & 0.0 \\
50\,000  & 0.0 \\
20\,000  & 0.0 \\
8\,000   & 0.0 \\
4\,000   & 0.0 \\
2\,000   & 0.0 \\
1\,000   & 0.0 \\
\bottomrule
\end{tabular}
\end{center}

\noindent
Exactly zero, not a tolerance. \code{examples/invariants} adds twelve pinned
witnesses across five tape shapes, each run whole and chunked: the
\code{lattice} case gives an identical Jacobian hash whole, at 5 chunks, at 16
chunks, and under both elimination directions.

\section{What replay costs}
\label{sec:replaycost}

Measured on a synthetic tape of 16 variables and 120 steps
($\approx$1\,920 recorded operations), GCC 13.3, \code{-O2}:

\begin{center}
\begin{tabular}{@{}lrrrr@{}}
\toprule
budget & partitions & time & vs.\ whole & replay share \\
\midrule
the program itself, no AD & --- & 0.0085\,ms & --- & --- \\
whole tape   & 1   & 5.78\,ms  & $1.00\times$ & --- \\
200\,kB      & 8   & 15.40\,ms & $2.66\times$ & 62\% \\
50\,kB       & 29  & 18.35\,ms & $3.18\times$ & 68\% \\
12\,kB       & 121 & 36.40\,ms & $6.30\times$ & 84\% \\
\bottomrule
\end{tabular}
\end{center}

Two conclusions. First, the differentiation machinery costs roughly $680\times$
the underlying arithmetic, so the target for any optimisation is the tape, never
the user's function. Second, the elimination cost is \emph{identical} at every
budget (33\,647 multiply--adds), so all of the growth in that table is replay.
Chunking is not free, and the cheapest speedup available to a user is a larger
budget.

\chapter{The memory budget}
\label{ch:mem}

\section{What a graph costs}

The budget is spent against a model of the graph's payload:

\begin{lstlisting}
static const largeint VERTEX_BYTES = sizeof(Vertex);                    // 72
static const largeint EDGE_BYTES   = sizeof(Edge) + 2*sizeof(AdjEntry); // 24 + 32 = 56
\end{lstlisting}

used at exactly three sites: \code{check\_memory()}, \code{get\_memory()}, and
the profiling branch of \code{Process::finalize()}.

\section{\svegp{24}: the accounting never counted adjacency}

The original model was \code{intmed\_count*sizeof(Vertex) + edge\_count*sizeof(Edge)}.
But an edge is not only its \code{Edge} object: it is also the \emph{two
adjacency entries that reference it}, one in the target's in-list and one in the
source's out-list. Before the arenas those were a pair of red-black tree nodes
and went uncounted; after the representation change they became a pair of
\code{AdjEntry} and were \emph{still} uncounted. Same omission, different
container. An edge costs $24+2\times16=56$ bytes of payload, not 24, so the
budget had been under-charging every edge by more than half since the beginning.

The constants deliberately exclude allocator overhead, container slack and arena
block granularity. Those are properties of the allocator, not of the graph, and
--- decisively --- they are not knowable during the profiling pass, which is
when the budget decision is made and when nothing has been allocated yet. So the
model is the graph's payload exactly, and the process always holds somewhat
more.

\section{The two errors were cancelling}

Correcting the edge charge moved the chunk counts \emph{back to where they had
started}:

\begin{center}
\begin{tabular}{@{}lrrl@{}}
\toprule
 & \code{sizeof(Vertex)} & per edge & lattice chunk counts \\
\midrule
original                  & 120 & 24 & 5, 16, 4 \\
after the representation work & 72 & 24 & 4, 11, 3 \\
after \svegp{24}          & 72 & \textbf{56} & \textbf{5, 16, 4} \\
\bottomrule
\end{tabular}
\end{center}

The original formula over-charged for vertices (120 against a real payload of
72) and charged nothing for adjacency, and the two errors roughly cancelled.
That is why the examples' hand-calibrated budgets worked at all, and why
recalibrating them turned out to be almost a no-op: \code{bratu} $14\to13$,
\code{hybrdj1} $20\to18$, \code{heat} $100\to99$.

This is the most instructive defect in the library's history. A wrong model
that happens to produce right-looking answers is harder to find than one that
produces wrong answers, and it is found only by measuring the thing the model
claims to predict.

\section{The drift alarm}

Because the model deliberately under-reports, the gap between it and reality
must be watched rather than assumed. \code{examples/invariants} measures process
heap against the accounted figure on two fixed shapes and fails outside a band:

\begin{lstlisting}
const double LOW = 1.2, HIGH = 1.9;
...
const double ratio = double(peak)/double(acct);
const bool ok = (ratio>=LOW && ratio<=HIGH);
\end{lstlisting}

The band is set by experiment, not by guesswork, and it has been re-pinned twice
as the representation changed:

\begin{center}
\begin{tabular}{@{}lll@{}}
\toprule
 & correct accounting & \svegp{24} reverted \\
\midrule
before the adjacency arena & 1.67 -- 1.70 & 2.38 -- 2.42 \\
after the adjacency arena  & \textbf{1.51 -- 1.56} & 2.16 -- 2.22 \\
\bottomrule
\end{tabular}
\end{center}

The first attempt at a band was $[1.0,3.0]$, and the control --- reverting
\svegp{24} --- walked straight through it; only the pinned chunk column caught
the reverted fix. The band was tightened to $[1.2,2.2]$. Then the adjacency
arena lowered the real heap, which lowered \emph{both} columns, and $2.16$
landed back inside $[1.2,2.2]$: the same failure, a second time. It is
$[1.2,1.9]$ now, verified in both directions.

The alarm is skipped under AddressSanitizer, which replaces the allocator
wholesale and makes \code{mallinfo2} report ASan's redzoned bookkeeping rather
than the program's footprint. The first ASan run of the alarm reported
``0 accounted, 0 heap, \code{-nan}x''.

\section{\svegp{25}: a budget that was not a budget}

\code{examples/heat} was the only example not taking a plain byte budget. Its
third argument was a \emph{multiplier}: the budget was \code{it*431208}, and
431\,208 was one time step's tape footprint measured at $n_x=400$. So \code{it}
meant ``how many $n_x{=}400$ time steps fit in a chunk'', and at any other grid
it silently meant something else --- at $n_x=40$ a step costs about a tenth as
much, so \code{it=1} bought ten steps per chunk rather than one.

It takes bytes now, and the knob finally behaves like one:

\begin{center}
\begin{tabular}{@{}lr@{}}
\toprule
invocation & chunks \\
\midrule
\code{40 100 400000}  & 11 \\
\code{40 200 400000}  & 22 \hspace{1em}\emph{double the steps, double the chunks} \\
\code{100 400 400000} & 107 \\
\code{100 400 800000} & 54 \hspace{1em}\emph{double the budget, half the chunks} \\
\bottomrule
\end{tabular}
\end{center}

\noindent
The change had a hazard: \code{main.cpp} and \code{RUNARGS} in the example's
Makefile had to agree, and bytes in one with a stale \code{1} in the other is a
one-byte budget, which chunks on every operation and reads as a hang rather than
a mistake. So \code{main.cpp} rejects an implausibly small budget and names the
exact fix. A stale Makefile now fails in a second with instructions instead of
appearing to hang.

\chapter{The tape object}
\label{ch:tape}

\section{\svegp{26}: two global states that had to agree}

The library used to have not one global but two, and they had to be kept
consistent by hand: \code{Process::instance}, a singleton holding the graph;
and about a dozen file-statics in \code{API.cpp} holding everything else about
the recording --- \code{run\_target}, \code{run\_counter},
\code{independent\_size}, the 2-D dimensions, and the two shadow-copy buffers.
Neither knew about the other.

\svegp{17} --- a stride bug that wrote past the end of the heap on any
non-square grid --- lived in the second one, in code that had no business
existing apart from the tape it described:

\begin{lstlisting}
// the row stride of a dep_y_dim x dep_x_dim grid is dep_x_dim (the COLUMNS),
// not dep_y_dim.  Flattening with dep_y_dim aliased entries onto each other
// and ran off the end of dep_shadow_copy as soon as the grid was not square.
for( largeint i=0 ; i<tp->dep_y_dim ; i++ ){
  for( largeint j=0 ; j<tp->dep_x_dim ; j++ ){
    tp->dep_shadow_copy[tp->dep_x_dim*i+j] = y[i][j].val;
  }
}
\end{lstlisting}

Every example used a square grid --- \code{bratu} calls
\code{set\_indep\_dimension(n,n)} --- which is why nothing caught it for years.

Both states are now members of one object:

\begin{lstlisting}
class TapeState
{
public:
  Process proc;                    // the graph, by value

  largeint independent_size, dependent_size, memory_available;
  largeint indep_x_dim, indep_y_dim, dep_x_dim, dep_y_dim;
  largeint run_target, run_counter;

  double * indep_shadow_copy;
  double * dep_shadow_copy;
  ...
private:
  TapeState( const TapeState & );          // owns buffers and a graph
  TapeState & operator=( const TapeState & );
};
\end{lstlisting}

and the only mutable object left at namespace scope in the entire library is one
pointer:

\begin{lstlisting}
static thread_local TapeState * g_current_tape = NULL;
\end{lstlisting}

\noindent
\code{nm} on the archive confirms it: one entry in the data segment. Everything
else static is \code{const}.

\paragraph{No public API changed.} The operators still \emph{find} the tape
rather than being handed one, so no signature moved and no user program was
touched --- which is precisely what let the existing verification net check the
whole change. All twelve pinned invariants came out bit-identical.

\section{What that bought, precisely}

Being exact about which limitation came from which decision was the point of the
exercise:

\begin{itemize}
\item \textbf{Two tapes at once} --- structurally impossible under the
      singleton, since \code{get\_proc\_instance()} handed the same
      \code{Process} to everybody. Now routine.
\item \textbf{Threads} --- the \code{thread\_local} pointer removes the
      structural reason the library could never be thread-safe.
      Chapter~\ref{ch:threads} measures what that actually delivers.
\item \textbf{Not nesting.} A tape whose edge weights are themselves
      \code{active} --- what an exact Hessian needs --- requires the scalar type
      to be templated. That is a separate and larger change, and removing the
      singleton does not begin it.
\end{itemize}

\section{\svegp{27}: the compiler found what had to move, not what had to be guarded}

Deleting the file-statics and letting the build fail was the right way to
enumerate the work: 29 uses of \code{myProc} and about 40 of the rest, all named
precisely. But the technique has a blind spot, and it bit.

A function reading a \emph{file-static} was reading something that always
existed and was simply zero. The same function reading \code{tp->field}
dereferences a pointer that is null when no tape is open. The code compiles
either way.

\code{checkpoint()} was the sharp one. With no tape it used to return
\code{true} on the first call --- \code{run\_counter} was 0 --- run the caller's
section with nothing recording, then null-dereference on the \emph{second} call.
Seven entry points needed guards no compiler was going to ask for. They are
guarded now, preserving the old tolerant answers where those were sensible
(\code{get\_partitions()} still returns 0 with no tape), and \code{checkpoint()}
ends the loop before it starts rather than running a section that records
nothing.

\code{examples/regress} now calls \emph{every} public entry point with no tape
open and asserts none of them crash. That is the test that would have caught it,
and it did not exist until the bug made the case for it.

\section{\svegp{27}: ownership belongs in a destructor}

A second, quieter defect followed from the same change. \code{initialize()} on
an already-open tape diagnoses the misuse and deletes the old tape --- but
\code{\textasciitilde TapeState} freed only the two shadow buffers, and
\code{Process} had \emph{no destructor at all}. The graph teardown lived only in
the productive branch of \code{Process::finalize()}, which that path never
calls. Every \code{Vertex} leaked, in proportion to the tape.

Reproduced under LeakSanitizer on a two-line program: 280 bytes in 6
allocations, with stacks in \code{vertex\_on\_lhs} and \code{vertex\_on\_rhs}.
It scales with the graph, so a program that re-initialises in a loop leaks the
whole thing each time.

The fix moved the teardown into a private \code{destroy\_graph()} that both
\code{\textasciitilde Process()} and \code{finalize()} call:

\begin{lstlisting}
void Process::destroy_graph()
{
  std::set<Vertex*> owned;

  for( ... intmed_map ... ) if(it->second) owned.insert(it->second);
  for( ... intmed_vec ... ) if(*it)        owned.insert(*it);

  for( std::set<Vertex*>::iterator it=owned.begin() ; it!=owned.end() ; it++ )
    retire(*it);

  intmed_map.clear();
  intmed_vec.clear();
  dep_vec.clear();

  edge_arena.clear();
  adj_arena.clear();
}
\end{lstlisting}

The \code{std::set} is deliberate. \code{intmed\_vec} and \code{intmed\_map}
never hold the same vertex at a moment a caller can observe --- each transfer
clears the source before returning --- but a teardown is the wrong place to lean
on that, and a graph abandoned mid-section is exactly where the two containers
are least likely to be in the state the happy path leaves them in.
\code{dep\_vec} is cleared rather than deleted from: every vertex in it is also
in \code{intmed\_map}, so deleting from both would be a double free.

\chapter{The RAII surface}
\label{ch:raii}

\section{Why}

The free-function protocol is eleven ordered steps over hidden global state, and
--- this is the part that matters --- four of them go wrong by being
\emph{forgotten} rather than by being written incorrectly: the
\code{try}/\code{catch}, \code{free\_jacobian}, \code{finalize}, and the
\code{delete[]}s. Every misuse the audit found was a silent memory error before
it was made a diagnostic. Diagnostics are second best. The best is making the
misuse unrepresentable:

\begin{lstlisting}
Tape t(n,m,budget);                     // ctor/dtor replace initialize/finalize
active * x = t.independents(x0);        // hands back the only thing you can use
t.run( y , [&]{ y = f(x); } );          // owns the loop AND the try/catch
t.dependents(y);
Jacobian J = t.harvest();               // a value; no free_jacobian
\end{lstlisting}

Each step returns or requires the thing the next one needs, so a step cannot be
skipped or reordered without the compiler saying so.

\section{A wrapper, not a second implementation}

This is the central implementation decision of the surface. Every member makes
its tape current for the duration and then calls the same free function the
examples have always called:

\begin{lstlisting}
maxwell::Jacobian maxwell::Tape::harvest( bool print_out )
{
  Scope s(*this);

  const largeint m = st_->dependent_size;
  const largeint n = st_->independent_size;

  double ** A = NULL;
  maxwell::harvest(m,n,A,print_out);

  if(!A) return Jacobian();

  Jacobian J(m,n);
  for( largeint i=0 ; i<m ; i++ )
    for( largeint j=0 ; j<n ; j++ ) J(i,j) = A[i][j];

  maxwell::free_jacobian(m,A);
  return J;
}
\end{lstlisting}

Consequently the twelve pinned invariants, the 68 finite-difference checks and
the budget-invariance sweep all cover the new surface, because underneath it
\emph{is} the old one. A reimplementation would have needed its own net. The
claim was then tested directly by migrating the verification harness itself to
the object surface: all twelve rows came out bit-for-bit identical.

\section{Scopes}

A \code{Tape} is not current outside its own calls. That is what lets two
\code{Tape} objects coexist without either being closed, and what makes a throw
out of a caller's section leave the current-tape pointer where it found it:

\begin{lstlisting}
class Scope
{
public:
  explicit Scope( Tape & t ) : t_(t) { t_.enter_(); }
  ~Scope(){ t_.leave_(); }
private:
  Tape & t_;
  Scope( const Scope & );
  Scope & operator=( const Scope & );
};
\end{lstlisting}

\section{\svegp{29}: a bug only a constructed test could find}

\code{enter\_}/\code{leave\_} originally kept a single saved-tape slot. That is
correct for nesting two \emph{different} tapes, and wrong for a tape re-entered
from inside its own section --- which happens the moment a caller reads
\code{t.memory()} while recording, as the drift alarm does:

\begin{lstlisting}
void maxwell::Tape::enter_()
{
  if( st_->scope_depth==0 ){ st_->saved = internals::current_tape(); }
  st_->scope_depth++;
  internals::set_current_tape(st_);
}

void maxwell::Tape::leave_()
{
  if( st_->scope_depth ) st_->scope_depth--;
  if( st_->scope_depth ) return;          // still inside an outer scope
  internals::set_current_tape( st_->saved );
  st_->saved = NULL;
}
\end{lstlisting}

Without the depth counter the inner scope overwrites the real predecessor with
the tape itself, and leaving the \emph{outer} scope restores \code{NULL} rather
than the enclosing tape. Nothing reports it: an enclosing section simply carries
on recording onto no tape at all and harvests a silently wrong Jacobian.

What is instructive is that the example which \emph{contains} the re-entrant
call does not detect the bug --- it has no enclosing tape for the clobber to
destroy. Only a deliberately constructed case in the regression suite, with one
tape running inside another's section, fails. The migration created the
hazardous call; a test written specifically for it revealed the consequence.

\section{Mixing the two surfaces}

Since the free functions remain public, a program could call \code{finalize()}
from inside a \code{Tape} section --- which would delete the state the
\code{Tape} object still points at, a use-after-free the moment the block ends.
Both \code{initialize()} and \code{finalize()} refuse when the current tape is
object-owned:

\begin{lstlisting}
if( open && open->owned ){
  std::cerr << "maxwell::initialize: a maxwell::Tape object is recording -- "
               "not opening a second tape over it.  Use one surface or the "
               "other, not both.\n";
  return;
}
\end{lstlisting}

Removing that guard and running the suite under AddressSanitizer produces
\code{heap-use-after-free}, which is how the guard is known to be load-bearing.

\section{Why the free functions stayed}

The obvious follow-up to an object surface is to retire what it wraps. That was
considered and rejected on evidence: three things the object surface cannot
express, all of them exercised by the regression suite.

\begin{itemize}
\item \textbf{Stepping the loop.} \code{run()} owns the whole checkpoint loop,
      which is exactly what makes it safe --- the \code{try}/\code{catch} cannot
      be forgotten. The price is that a caller cannot interleave two tapes
      mid-replay, which the strongest available test of the two-tape capability
      does.
\item \textbf{Testing the lower layer.} The free functions are still public and
      are what every \code{Tape} call turns into. Tests of their misuse
      behaviour are not obsolete; they are the only thing holding it.
\item \textbf{Cross-checking.} The suite compares the object path against the
      free-function path bit-for-bit. Migrating the reference would make that
      comparison compare the wrapper with itself.
\end{itemize}

So the answer is that they stay, as a documented lower layer.

\chapter{Storage: the two arenas}
\label{ch:arena}

\section{Where the memory went}

The original graph representation was two \code{std::map<largeint,Edge*>} per
vertex. That cost 96 of the vertex's 120 bytes before a single edge existed, and
every edge bought \emph{three} allocations: one for the \code{Edge} itself and
two red-black tree nodes, one in each direction, each with the allocator's own
header, scattered across the heap.

Measured on one \code{heat\_assim} tape ($n_x{=}20$, $n_t{=}40$):

\begin{center}
\begin{tabular}{@{}lr@{}}
\toprule
what \code{get\_memory()} accounted & 728\,856 bytes \\
what was actually on the heap       & 1\,780\,656 bytes \\
\textbf{ratio}                      & \textbf{2.44$\times$} \\
\bottomrule
\end{tabular}
\end{center}

\noindent
So the memory budget --- the feature the whole checkpointing design exists to
serve --- was wrong by a factor of two and a half. That is also why the budgets
in the examples were magic numbers found by experiment: the model did not
predict reality, so someone measured instead.

The representation was replaced in three steps, each verified against the pinned
invariants before the next began.

\section{Step 1: the edge arena}

Edges come from blocks and freed slots go on a free list, with the adjacency
maps left exactly as they were. The smallest change that removes one of the
three allocations per edge.

\begin{lstlisting}
Edge * EdgeArena::acquire( Vertex * src , Vertex * tgt , double eval )
{
  Edge * e = 0;

  if( !free_list.empty() ){
    e = free_list.back();
    free_list.pop_back();
  }else{
    if( blocks.empty() || block_used==BLOCK ){
      blocks.push_back( new Edge[BLOCK] );
      slots += BLOCK;
      block_used = 0;
    }
    e = blocks.back() + block_used;
    block_used++;
  }

  e->eval = eval;  e->src = src;  e->tgt = tgt;
  live++;
  return e;
}
\end{lstlisting}

Recycling was already the behaviour: \code{eliminate()} releases an in-edge and
then allocates fill-in inside the same loop, so the allocator could always have
handed back the block it had just freed. The free list makes that explicit
rather than introducing it --- which is why the pinned Jacobians did not move.

\begin{center}
\begin{tabular}{@{}lrr@{}}
\toprule
 & before & after \\
\midrule
allocations per tape & 39\,462 & \textbf{27\,908} \\
peak heap            & 1\,780\,720 B & \textbf{1\,559\,216 B} \\
under-count ratio    & 2.44$\times$ & \textbf{2.14$\times$} \\
\bottomrule
\end{tabular}
\end{center}

\noindent
The 11\,554 allocations that disappeared are exactly the per-edge \code{new
Edge}. What remained --- about 23\,000 tree nodes, two per edge --- confirmed the
diagnosis: \emph{the map nodes were the bulk, not the \code{Edge} objects}.

\section{Step 2: sorted adjacency}

The two \code{std::map}s became one sorted array per vertex per direction. The
choice of a sorted array over a linked list was made on the degree
measurements: the direct-link probe scans a predecessor's out-list of 40 to 109
entries, growing with problem size, so the search had to stay logarithmic.

\begin{center}
\begin{tabular}{@{}lrrr@{}}
\toprule
 & original & after step 1 & \textbf{after step 2} \\
\midrule
allocations per tape & 39\,462 & 27\,908 & \textbf{17\,941} \\
peak heap            & 1\,780\,720 & 1\,559\,216 & \textbf{820\,896} \\
\code{sizeof(Vertex)} & 120 & 120 & \textbf{72} \\
under-count ratio     & 2.44$\times$ & 2.14$\times$ & \textbf{1.63$\times$} \\
\bottomrule
\end{tabular}
\end{center}

\noindent
Peak memory down 54\% and allocations down 55\% from the start. All twelve
\emph{structural} pinned values --- every Jacobian hash and every elimination
cost --- were byte-for-byte identical. The \emph{accounting} column moved exactly
as predicted, $5\to4$, $16\to11$, $4\to3$ chunks, because \code{sizeof(Vertex)}
fell and the same byte budget therefore accounted for more graph. That split
between structural and accounting values is what let a single test report
``the derivative is unchanged and the partitioning moved'' rather than simply
``failed''.

\section{Step 3: the adjacency arena}

The last per-vertex allocation. A \code{std::vector} adjacency was one
allocation per vertex per direction the moment it became non-empty, plus a
reallocation on every doubling, plus the doubling's slack --- and elimination
creates and destroys vertices in the millions.

This step had to wait for the tape object, and the reason is instructive. An
arena needs somewhere to put the \emph{release}. A \code{Vertex} has no arena and
giving it a back-pointer would have cost eight of the seventy-two bytes the
arena exists to save. So \code{Process} owns the arena and is the only thing
that destroys a vertex:

\begin{lstlisting}
void Process::retire( Vertex * v )
{
  if(!v) return;
  v->in_edges.clear(adj_arena);
  v->out_edges.clear(adj_arena);
  delete v;
}
\end{lstlisting}

\noindent
and \code{\textasciitilde Vertex()} is deliberately empty, with a comment saying
so, because an empty destructor that looks like an oversight is worse than no
destructor at all.

The arena carves fixed size classes from slabs and recycles them per class:

\begin{lstlisting}
AdjEntry * AdjArena::acquire( largeint want , largeint & cap )
{
  const largeint c = class_of(want);
  cap = cap_of(c);
  ...
  if( !recycled[c].empty() ){
    AdjEntry * p = recycled[c].back();
    recycled[c].pop_back();
    return p;
  }

  if( left[c] < cap ){
    largeint slab = next_slab[c] ? next_slab[c] : SLAB_MIN;
    if( slab < cap ) slab = cap;
    AdjEntry * p = new AdjEntry[size_t(slab)];
    slabs.push_back(p);
    entries += slab;  cursor[c] = p;  left[c] = slab;
    next_slab[c] = (slab<SLAB_MAX) ? slab*2 : SLAB_MAX;
  }

  AdjEntry * p = cursor[c];
  cursor[c] += cap;  left[c] -= cap;
  return p;
}
\end{lstlisting}

\subsection*{Two parameter choices, both measured wrong first}

This is worth recording in full, because both wrong choices looked right.

\paragraph{Flat 4\,096-entry slabs.} The first attempt. It hit the target ---
allocations fell 73\% --- and pushed \emph{peak heap up} by 20 to 32\%. A size
class holding three blocks paid for all 64\,kB of its slab. Replaced with
geometric growth per class.

\paragraph{Geometric growth barely helped}, which was the useful surprise: the
real cost was the \emph{minimum block size}. Classes started at four entries.
But most vertices on a tape have in-degree one or two and out-degree one --- an
elementary operation has one or two operands --- so the commonest allocation in
the library was overshooting by two to four times, on the two blocks that
\emph{every} vertex carries. Capacities start at one entry now, and that is
where the win came from.

Neither would have been found by reasoning about it. The first looked correct
and was a regression; the second was invisible until the first was fixed.

\subsection*{Result}

\begin{center}
\begin{tabular}{@{}lrr@{}}
\toprule
shape & allocations & peak heap \\
\midrule
12 vars $\times$ 40 steps & 9\,297 $\to$ \textbf{2\,564} ($-72\%$) & 425\,680 $\to$ \textbf{362\,640} ($-15\%$) \\
20 vars $\times$ 25 steps & 9\,766 $\to$ \textbf{2\,707} ($-72\%$) & 639\,040 $\to$ \textbf{595\,360} ($-7\%$) \\
20 vars $\times$ 40 steps & 15\,467 $\to$ \textbf{4\,213} ($-73\%$) & 876\,896 $\to$ \textbf{811\,568} ($-7\%$) \\
\bottomrule
\end{tabular}
\end{center}

\noindent
\code{sizeof(Vertex)} is still 72 --- \code{Adjacency} is a pointer and two
counters, the same 24 bytes a \code{std::vector} occupied --- so
\code{VERTEX\_BYTES} is unchanged and \emph{no chunk count moved}.

The effect on wall time was larger than the memory figures suggest and was not
anticipated. Reproducibly, across three runs:

\begin{center}
\begin{tabular}{@{}lrr@{}}
\toprule
 & before & after \\
\midrule
whole tape          & 13.8 -- 15.4\,ms & \textbf{5.4 -- 5.7\,ms} \\
chunked, 200\,kB    & 22.3 -- 23.5\,ms & \textbf{14.7 -- 15.2\,ms} \\
\bottomrule
\end{tabular}
\end{center}

\noindent
That is $2.5\times$ on an allocation-heavy tape. It is a reminder that
allocation count, not resident bytes, was the dominant cost all along.

\chapter{Concurrency}
\label{ch:threads}

\section{What is shared}

One object at namespace scope in the whole library is mutable, and it is
\code{thread\_local}. Everything else static is \code{const}: two size
constants, one mathematical constant, and the arena class parameters. Every
piece of state a recording touches --- the graph, both arenas, the counters, the
shadow buffers, the 2-D dimensions --- lives in the \code{TapeState} that
\svegp{26} collected them into.

\section{What works}

Independent tapes on separate threads. Verified under ThreadSanitizer:

\begin{itemize}
\item 4 threads $\times$ 25 tapes each, half of them chunking hard so
      \code{BreakException} unwinds on several threads simultaneously: no races,
      every derivative correct.
\item 8 threads $\times$ 6 rounds, each round building a deep tape three ways
      --- object surface whole, object surface chunked at 3\,000 bytes, and the
      free functions at 4\,000 --- and cross-checking all three: zero
      disagreements, no races. The free functions are as thread-safe as the
      object surface, because they go through the same thread-local pointer.
\item \textbf{Negative control:} one \code{Tape} object driven from two threads
      races immediately, reported in \code{Tape::enter\_()} and
      \code{independents()}. The boundary is where it should be, and the
      sanitizer says so loudly rather than corrupting quietly.
\end{itemize}

Measured speedup over 16 independent tapes on 2 cores:

\begin{center}
\begin{tabular}{@{}lrrr@{}}
\toprule
budget & serial & parallel & speedup \\
\midrule
whole tape & 90.6\,ms & 50.7\,ms & $1.79\times$ \\
200\,kB    & 247.3\,ms & 129.0\,ms & $1.92\times$ \\
50\,kB     & 308.4\,ms & 157.4\,ms & $1.96\times$ \\
12\,kB     & 561.7\,ms & 281.1\,ms & $2.00\times$ \\
\bottomrule
\end{tabular}
\end{center}

\noindent
Efficiency improves slightly at tighter budgets, because replay is pure
per-thread arithmetic while recording contends on the allocator.

The memory trade is milder than it first appears. With a fixed 200\,kB pool: one
thread at 200\,kB is 15.4\,ms per tape; two threads at 100\,kB each is about
16.4\,ms per tape, two at a time, for an effective 8.2\,ms --- $1.88\times$.
Roughly 94\% of two cores survives halving the budget.

\section{What does not work, and why}

Parallelising \emph{across partitions} is not available, and the obstacle is
architectural rather than mechanical. Three reasons, worst first.

\paragraph{Value dependency.} Partition $k$'s inputs are the program state after
partitions $1..k-1$ have executed. The replay design exists precisely because
that state is never stored. Building partitions concurrently means checkpointing
the \emph{values} at boundaries --- which is the ``spill the tape rather than
replay the program'' alternative, a different design with different costs, not a
parallelisation of this one.

\paragraph{Elimination order.} Each pass folds its partition into the
accumulating result, last to first, each step consuming what came before.
Chapter~\ref{ch:elim} showed that elimination order changes the result in the
last bits. Reordering or interleaving partitions therefore changes the Jacobian,
destroying the bit-exactness across budgets that has twelve pinned witnesses ---
which is the property that makes chunking trustworthy in the first place.

\paragraph{Shared recording state.} \code{next\_vertex\_idx},
\code{intmed\_vec}, the partition-boundary logic and both arenas are per-tape,
not per-thread. Lockable --- but locking makes index assignment
nondeterministic, which returns to the previous paragraph.

The same three apply to placing an OpenMP loop \emph{inside} a section. Making
that work would need per-thread sub-tapes with index ranges reserved during the
profiling pass --- which already counts everything, so it could hand out
deterministic ranges --- followed by a merge. That is a substantial design, and
it only pays if the section contains genuinely independent work, which a deep
time-stepping tape does not.

\paragraph{One avenue remains open.} Recording partition $k+1$ and eliminating
partition $k$ are order-independent: recording needs values, not the elimination
result. A two-stage pipeline would preserve the accumulation chain. Its ceiling
is the smaller of the two phases, which has not been measured, and it would
require the arenas to be split per stage. It is noted here as the one form of
intra-tape parallelism the architecture does not forbid.

\section{A caveat on diagnostics}

Diagnostics go to \code{std::cerr} unsynchronised. That is not a data race ---
the standard stream objects are safe --- but two threads misusing the API
simultaneously will interleave their messages mid-line.

\chapter{Verification}
\label{ch:verify}

The way this library is tested is itself a design choice, and the report would
be incomplete without it.

\section{Structural values and accounting values}

\code{examples/invariants} freezes what the library computes across twelve cases
and five tape shapes, and pins \emph{two kinds} of number:

\begin{itemize}
\item \textbf{Structural} --- the elimination cost, and a bit-exact FNV-1a hash
      of every harvested Jacobian entry. These are properties of the graph and
      the arithmetic. They must never move.
\item \textbf{Accounting} --- the partition count. This is \emph{expected} to
      move when the memory model changes, and is pinned so that the move is
      deliberate.
\end{itemize}

The split earns its keep every time the representation changes. A single-column
test would have reported a failure at step 2 of the arena work and told you
nothing about whether the derivative had changed.

The hash is a refactoring change-detector, not a portable golden value, so the
example compiles with \code{-ffp-contract=off}: otherwise the compiler may fuse
\code{a*b+c} into an FMA depending on the optimisation level, and the hash would
track \code{-O} rather than the library.

\section{Controls}

\emph{A test that has never failed is not known to work.} Every assertion added
during this work was checked by deliberately breaking the thing it guards:

\begin{center}
\begin{longtable}{@{}p{7.2cm}p{7.2cm}@{}}
\toprule
control & observed \\
\midrule
\endhead
reverse the out-edge iteration order & nothing --- correctly \\
drop fill-in edges below $10^{-12}$ & 1 case, on \emph{cost} \\
perturb every partial by 1 ulp & all 12, on \emph{hash} \\
revert \svegp{24} (uncount adjacency) & chunk column, and the drift band \\
\code{\textasciitilde Process()} emptied & LeakSanitizer, 560 bytes in 12 allocations \\
\code{run()} stops catching \code{BreakException} & abort, exit 134 \\
object-owned guard removed from \code{finalize()} & ASan \code{heap-use-after-free} \\
\code{leave\_()} stops restoring & the nested-tape case fails \\
\code{set\_current\_tape} made a singleton again & all three two-tape assertions fail \\
off-by-one in the arena's block capacity & ASan \code{heap-use-after-free} \\
\bottomrule
\end{longtable}
\end{center}

\noindent
Two entries in that table are the most valuable. The first --- reversing
adjacency order and observing \emph{nothing} --- is a finding, not a null
result: it disproved a constraint the design had assumed and freed the storage
layout. And one control that was run \emph{did not} fail: removing the
current-tape restore from \code{\textasciitilde Tape()} breaks nothing, because
every scope pairs its enter with its leave even while an exception unwinds. That
line stays as belt-and-braces, and the code comment says plainly that the suite
does not cover it. Recording a control that passed is as useful as recording one
that failed.

\section{Sanitizers}

\code{make sanitize} rebuilds into a \emph{separate} tree under AddressSanitizer
and LeakSanitizer. The separate tree is deliberate: \code{make} compares
timestamps only, so sanitizing in place leaves every binary newer than its
source and a following plain \code{make} rebuilds nothing --- you keep running
instrumented code without being told, which is how a fast example comes to look
like a hang.

\section{What the suite asserts}

\begin{center}
\begin{tabular}{@{}lp{9.2cm}@{}}
\toprule
example & what it holds \\
\midrule
\code{invariants} & 12 pinned structural values; the accounting drift band \\
\code{fdcheck}    & 68 finite-difference checks of every elementary function \\
\code{regress}    & every fixed defect, both API surfaces, budget invariance \\
\code{bratu}      & the 2-D checkpoint overload; asserts chunking \\
\code{heat}       & a 107-partition tape; asserts chunking \\
\code{heat\_assim}& adjoint vs finite differences; elimination direction; CG \\
\code{hybrdj1}    & a full $20\times20$ Jacobian; asserts chunking \\
\code{pinn}       & 2\,502 tapes; chunked gradient equals unchunked reference \\
\code{test\_xor}  & 4\,001 tapes; AD gradient vs central differences \\
\bottomrule
\end{tabular}
\end{center}

\part{Using the Library}

\chapter{Getting started}

\section{Building}

\begin{lstlisting}[language=bash]
make            # the library and every example, into build/
make test       # build, then run the whole suite in order
make list       # the examples and the arguments they run with
make sanitize   # rebuild into build-asan/ under ASan+LSan, run the suites
make clean      # remove build/ and build-asan/
\end{lstlisting}

Everything generated lands under \code{build/}: \code{build/lib/libmaxwell.a},
\code{build/obj/}, and one binary per example in \code{build/bin/}. Nothing is
written into the source tree.

A single example builds standalone from a cold tree --- \code{cd examples/pinn
\&\& make} recurses into \code{src/} first --- and \code{cd src \&\& make} builds
the library on its own.

\section{Including it}

\begin{lstlisting}
#include "maxwell.hpp"

using namespace maxwell;   // or qualify: maxwell::active, maxwell::sin
\end{lstlisting}

The header does not open the namespace for you. It brings in the \code{active}
scalar, the free-function API, \code{Tape} and \code{Jacobian}, and
\code{BreakException} --- and deliberately nothing internal, so the graph
representation can change underneath you without your code changing. It has,
twice.

\section{The shape of a program}

\begin{lstlisting}
#include "maxwell.hpp"
using namespace maxwell;

int main()
{
  const double x0[2] = { 1.1 , 2.3 };

  Tape t( 2 , 1 , 1u<<20 );          // 2 independents, 1 dependent, 1 MB budget

  active * x = t.independents(x0);   // the tape's own, values set and registered
  active y;

  t.run( y , [&]{ y = sin(x[0]*x[1]); } );

  t.dependent(y);

  Jacobian J = t.harvest();

  std::printf("dy/dx0 = %.10f\n", J(0,0));
  return 0;
}
\end{lstlisting}

Five calls. The constructor opens the tape and the closing brace closes it, on
every path out of the block including a \code{return} or a \code{throw}.

\chapter{The tape surface}

\section{Construction}

\begin{lstlisting}
Tape t( n , m , memory_budget );
\end{lstlisting}

$n$ independents, $m$ dependents, and a budget \emph{in bytes}. The budget is
what the library may hold as graph payload; see Section~\ref{sec:budget} for how
to choose it. A \code{Tape} is not copyable --- it owns a graph.

\section{Independents}
\label{sec:independents}

\begin{lstlisting}
active *  independents( const double * x0 );
active ** independents( largeint rows , largeint cols , const double * x0 );
\end{lstlisting}

The tape hands back its own independents, with the values already in them and
already registered. This replaces three steps that had to agree with each other
--- declare, set values, register in the same order --- with one, and it closes
the bounds error that \svegp{18} was about.

The storage belongs to the tape and is sized exactly once, so the reallocating
\code{std::vector<active>} hazard of Section~\ref{sec:pitfalls} cannot happen to
it. The returned pointer is valid until the \code{Tape} dies; do not
\code{delete} it.

The 2-D form takes \code{rows} and \code{cols} in the order the indices are
written, \code{x[i][j]} with $i$ over rows, and reads \code{x0} row-major. The
underlying \code{set\_indep\_dimension} takes them the other way round, columns
first; that transposition is \svegp{17}, it is invisible on a square grid, and
it is now made inside the library where a caller cannot get it wrong.

\section{Running the section}

\begin{lstlisting}
template<class Body> void run( active  & y , Body body );   // one dependent
template<class Body> void run( active  * y , Body body );   // m dependents
template<class Body> void run( active ** y , Body body );   // a rows x cols grid
\end{lstlisting}

\code{run} drives the checkpoint loop to completion and owns the
\code{try}/\code{catch}. Call it after \code{independents()} and before
\code{dependents()}. For the 2-D form, call \code{dependent\_shape(rows,cols)}
first --- \code{checkpoint} reads the dependent shape on its first pass, so it
cannot be inferred later.

\section{Dependents, harvest, and the rest}

\begin{lstlisting}
void dependent ( const active & y );      // one
void dependents( const active * y );      // m of them
void dependents( active ** y );           // the dependent_shape() grid

Jacobian harvest( bool print_out = false );

largeint partitions();          // how many chunks; 1 means it never chunked
largeint cost();                // elimination cost, in multiply-adds
largeint memory();              // graph payload in bytes, by the budget's model
void     set_elim_mode( elim_t mode );   // before run(); default REVERSE_ELIM
\end{lstlisting}

\section{The Jacobian}

\begin{lstlisting}
class Jacobian
{
public:
  largeint rows() const;
  largeint cols() const;
  bool     empty() const;                       // harvest() refused
  double   operator()( largeint i , largeint j ) const;
  double & operator()( largeint i , largeint j );
  ...
};
\end{lstlisting}

A value, stored as one flat row-major block. Copying one copies it; there is
nothing to free. \code{empty()} is true when \code{harvest()} refused --- a
shape mismatch, or a program that never drove \code{run()} --- and it is worth
checking, because it is the one way the call can decline.

\chapter{The checkpoint contract}

This chapter is short and is the most important one in Part~II.

\medskip
\noindent
The code inside \code{run()} is executed \textbf{$1+N$ times}, where $N$ is the
number of partitions and is not known until the first pass has finished. It may
also be \textbf{abandoned part-way} by an exception thrown from inside an
arithmetic operator. Your section must therefore be:

\paragraph{Re-runnable.} Same inputs, same result, every time. The library
restores the independents between passes; anything else your section reads must
not have changed.

\paragraph{Side-effect free.} The profiling pass runs your section an extra
time. Printing from inside a section prints $N+1$ times. Drawing a random number
inside a section draws a different one on each pass and silently breaks
re-runnability --- put the RNG \emph{outside}, as \code{examples/pinn} does.

\paragraph{Exception-safe.} \code{run()} owns the \code{try}/\code{catch}, so
you cannot forget it or catch the wrong thing --- but an exception still unwinds
through \emph{your} code. Anything held in a raw pointer between the start of
the section and the throw point is leaked. Use containers, or let the section
own nothing:

\begin{lstlisting}
// WRONG: leaks nx+1 actives on every partition boundary
active * old = new active[nx+1];
...
delete [] old;                       // never reached when BreakException fires

// RIGHT
std::vector<active> old(nx+1);       // sized once; destructor runs while unwinding
\end{lstlisting}

\noindent
That exact defect --- raw \code{new[]} buffers swapped by pointer, leaking on
every break, thousands of times across a conjugate-gradient run --- is a real one
from this library's history, and nothing would have caught it but a sanitizer.

\chapter{Choosing a budget}
\label{sec:budget}

\section{What the number means}

The budget is bytes of \emph{graph payload}: vertices at
\code{sizeof(Vertex)} = 72 and edges at \code{sizeof(Edge)} + 2\,$\times$\,
\code{sizeof(AdjEntry)} = 56. It deliberately excludes allocator overhead and
container slack, because those are not knowable during the profiling pass, when
the decision is made. Expect the process to hold roughly 1.5 to 1.9 times the
budget --- \code{examples/invariants} measures that ratio on every run and fails
if it drifts.

\section{How to pick one}

Start with the memory you can afford, divided by about 1.7. Then check what you
got:

\begin{lstlisting}
Jacobian J = t.harvest();
std::printf("tape broken into %lu chunk(s)\n", (unsigned long)t.partitions());
\end{lstlisting}

\code{partitions()==1} means the budget was never reached and the checkpoint
machinery never fired. That is the right answer for a small problem, and a
warning sign for a large one --- it means you are holding the whole tape.

\section{What chunking costs}

From Section~\ref{sec:replaycost}: 8 partitions cost $2.7\times$ the whole-tape
time, 29 cost $3.2\times$, 121 cost $6.3\times$. The elimination work is
identical in every case; the growth is entirely the replay of your section.

So the ordering of remedies for a slow differentiation is:

\begin{enumerate}
\item raise the budget if the memory exists --- this is by far the largest lever;
\item choose the elimination direction (Section~\ref{sec:direction});
\item run independent tapes on separate threads (Section~\ref{sec:parallel}),
      worth up to the core count, and slightly less because each tape then gets
      a smaller share of memory.
\end{enumerate}

\section{Elimination direction}
\label{sec:direction}

\begin{lstlisting}
Tape t(n,m,budget);
t.set_elim_mode(REVERSE_ELIM);      // the default
// t.set_elim_mode(FORWARD_ELIM);
\end{lstlisting}

Must be set before \code{run()}. Both directions produce the same Jacobian ---
to within the last bits, since they accumulate in different orders --- at very
different cost. Reverse wins on a deep narrow tape, and the gap grows with
depth: \code{examples/heat\_assim} measures $6\,835$ multiply--adds reverse
against $112\,102$ forward, a $16.4\times$ ratio, on the same graph. Forward can
win when there are few independents and many dependents. \code{cost()} reports
the figure, so the honest way to choose is to measure both once on a
representative problem.

\chapter{Use cases}

Six worked cases, each taken from a shipped example that runs under
\code{make test}.

\section{A gradient: many inputs, one output}

The commonest case in optimisation and machine learning. With one dependent,
the $1\times n$ Jacobian \emph{is} the gradient.

\begin{lstlisting}
static double gradient( const double * w , double * g , largeint mem )
{
  Tape t( largeint(NP) , 1 , mem );

  active * p = t.independents(w);
  active L;

  t.run( L , [&]{ L = mse(p); } );

  t.dependent(L);

  Jacobian J = t.harvest();
  if(J.empty()){ return 0.0; }

  for( int i=0 ; i<NP ; i++ ) g[i] = J(0,i);

  return passive_value(L);        // the loss, as a plain double
}
\end{lstlisting}

\noindent
\emph{From \code{examples/test\_xor}, which trains a 2--4--1 network on XOR for
4\,001 epochs, one tape per epoch, 8 partitions each.} Note what the signature
says: parameters go in as \code{double} and come back as \code{double}. They are
active only \emph{inside} the tape. That matters --- see
Section~\ref{sec:pitfalls}.

The gradient is checked against central differences on an independently written
plain-\code{double} forward pass, which is the only way such a check proves
anything: if both sides were the same code, agreement would be meaningless.

\section{A full Jacobian: $n$ inputs, $n$ outputs}

\begin{lstlisting}
Tape t( n , n , budget );

std::vector<double> x0(n,-1.0);
active * x = t.independents( &x0[0] );

active * fvec = new active[n];

t.run( fvec , [&]{ fcn(n,x,fvec); } );

t.dependents(fvec);

Jacobian J = t.harvest(true);     // true keeps the FJAC listing

const double want = 3.0 - 4.0*(-1.0);
for( int k=0 ; k<n ; k++ ) assert( J(k,k)==want );
\end{lstlisting}

\noindent
\emph{From \code{examples/hybrdj1}, MINPACK's \code{hybrj1} test function at
$n=20$, tape broken into 18 chunks.} The dependents here are the caller's own
array, not the tape's --- only the independents come from the tape.

\section{A two-dimensional grid}

When the natural indexing is \code{x[i][j]}, the tape can hand out a grid:

\begin{lstlisting}
std::vector<double> x0(n*n,0.0);
... fill x0 row-major, boundary conditions and all ...

Tape t( n*n , n*n , (largeint)bs );

active ** x = t.independents( n , n , &x0[0] );
t.dependent_shape( n , n );

t.run( x , [&]{ bratu(n,x,r); } );    // x is both independent and dependent
t.dependents(x);
\end{lstlisting}

\noindent
\emph{From \code{examples/bratu}, a 2-D Bratu discretisation on a $50\times50$
grid at a 400\,kB budget, 13 chunks.} The independents are also the dependents
here, which is legal and common for an iterated map.

This is the shape \svegp{17} lived in. Because \code{bratu} is square, it could
never have exposed the bug; the regression suite drives a deliberately
non-square $3\times2$ grid for that reason.

\section{An adjoint: data assimilation}

A deep tape --- the whole point of the library.

\begin{lstlisting}
static double gradient( const double * w , const double * obs , double * g ,
                        elim_t mode , largeint mem )
{
  Tape t( largeint(NP) , 1 , mem );
  t.set_elim_mode( mode );

  active * p = t.independents(w);
  active Jc;

  t.run( Jc , [&]{ Jc = cost_active(p,obs); } );

  t.dependent(Jc);

  const double value = passive_value(Jc);

  Jacobian J = t.harvest();
  for( int i=0 ; i<NP ; i++ ) g[i] = J(0,i);

  return value;
}
\end{lstlisting}

\noindent
\emph{From \code{examples/heat\_assim}: a 1-D heat equation marched 40 steps,
with the initial condition as the control and 36 observations as data. The cost
functional's gradient is the adjoint.} The example then drives linear conjugate
gradient with exact Hessian-vector products obtained by differencing the AD
gradient, and recovers the initial condition to an rms error of
$1.0\times10^{-13}$ over 19 controls, taking the cost from $7.71$ to
$3.6\times10^{-29}$.

The tape here is 40 time steps chained end to end and is eliminated whole rather
than per partition --- which is what makes the $16.4\times$ reverse-versus-forward
ratio visible.

\section{Verifying a derivative}

Never trust an AD result you have not checked once. The pattern:

\begin{lstlisting}
Tape t(1,1,1u<<30);
active * x = t.independents(&x0);
active y;
t.run( y , [&]{ y = f(x[0]); } );
t.dependent(y);
Jacobian J = t.harvest();
const double ad = J.empty() ? NAN : J(0,0);

const double h  = 1e-6;
const double fd = ( f_passive(x0+h) - f_passive(x0-h) ) / (2.0*h);

assert( std::fabs(ad-fd)/(1.0+std::fabs(fd)) < 2e-5 );
\end{lstlisting}

\noindent
\emph{From \code{examples/fdcheck}, which does this for every elementary function
the library overloads --- 68 checks, zero failures.} Use a generous tolerance:
you are checking the AD, and the finite difference is the inaccurate side.

\section{Independent tapes in parallel}
\label{sec:parallel}

Each thread builds its own tape. Nothing is shared, and no library call is
needed to arrange it:

\begin{lstlisting}
#pragma omp parallel for schedule(static)
for( int i=0 ; i<NCASE ; i++ ){
  Tape t( n , 1 , budget );
  active * x = t.independents(&starts[i*n]);
  active y;
  t.run( y , [&]{ y = f(x); } );
  t.dependent(y);
  Jacobian J = t.harvest();
  for( int j=0 ; j<n ; j++ ) out[i*n+j] = J(0,j);
}
\end{lstlisting}

Measured $1.79$ to $2.00\times$ on two cores, verified race-free under
ThreadSanitizer. Natural fits: finite-difference Hessians ($n$ independent
gradient evaluations), parameter sweeps, multi-start optimisation, ensembles,
and verification suites.

\paragraph{The one rule.} A \code{Tape} object belongs to one thread. Sharing
one across threads races immediately --- ThreadSanitizer reports it in
\code{Tape::enter\_()} --- and there is no locking mode that makes it safe,
because index assignment would become nondeterministic and the Jacobian would
stop being reproducible.

\paragraph{And the caveat.} $P$ parallel tapes need $P$ times the memory.
Checkpointing exists to bound memory, so this trades back some of what it buys.
The trade is mild --- halving the budget to double the threads keeps about 94\%
of the second core --- but it is real.

\chapter{Pitfalls}
\label{sec:pitfalls}

\section{\code{active} is not a value type}

Copying one records a vertex and an edge. Usually that is invisible. It stops
being invisible in a container that reallocates:

\begin{lstlisting}
// 5.5x the tape it should be: 15 reallocation copies, each recording
std::vector<active> h;
for( int i=0 ; i<8 ; i++ ) h.push_back( something(i) );

// fine: sized once, no element ever copied
std::vector<active> h(8);
for( int i=0 ; i<8 ; i++ ) h[i] = something(i);
\end{lstlisting}

\noindent
The derivative is correct either way. The tape is five and a half times bigger,
which changes how many partitions you get and how long it takes, and nothing
warns. \emph{Size your containers once, or reserve.}

\section{Do not hold actives across tapes}

An \code{active} that outlives its tape carries an index into a graph that has
been given back. Keep long-lived state --- optimiser parameters, iterates,
anything that survives a tape --- as \code{double}, and let
\code{t.independents()} make it active for the duration:

\begin{lstlisting}
// WRONG
active * p = new active[NP];
for( int e=0 ; e<epochs ; e++ ){ gradient(p,g,mem); for(...) p[i] = w[i]; }

// RIGHT
double * w = new double[NP];
for( int e=0 ; e<epochs ; e++ ){ gradient(w,g,mem); for(...) w[i] -= lr*g[i]; }
\end{lstlisting}

\noindent
Three shipped examples did the first thing and were changed to do the second.

\section{Do not mix the two surfaces}

Calling \code{finalize()} from inside a \code{Tape} section would free the state
the object still points at. The library refuses and says so, rather than letting
it happen --- but the diagnostic means your program is doing something it should
not.

\section{The budget is bytes}

Not kilobytes, not a multiplier, not a number of time steps. One example
historically took a multiplier and it silently meant a different thing at every
grid size; it takes bytes now, and rejects an implausibly small value rather
than chunking on every operation and appearing to hang.

\section{Check \code{empty()}}

\code{harvest()} returns an empty \code{Jacobian} rather than crashing when it
refuses --- a transposed shape, or a program that never drove \code{run()}. If
you index it without checking, you get zeros rather than a diagnostic.

\appendix

\chapter{Defect log}

Defects found and fixed during the 2026 audit, referred to throughout by
number. Numbering starts at 16; earlier numbers belong to the pre-audit history.

\begin{center}
\begin{longtable}{@{}lp{12.4cm}@{}}
\toprule
id & summary \\
\midrule
\endhead
\svegp{03} & harvest wrote the Jacobian column at the independent's vertex index rather than its position \\
\svegp{04} & harvest did nothing while still profiling, leaving the caller's pointer uninitialised \\
\svegp{05} & the requested Jacobian shape was trusted; a transposed request overran the heap \\
\svegp{06} & harvest printed unconditionally; \code{printf} with a non-literal format and \code{largeint} in \code{\%4d} \\
\svegp{07} & no way to release the harvested \code{double**}; every harvest leaked \\
\svegp{16} & \code{operator>=} declared but never defined --- a link failure \\
\svegp{17} & the 2-D checkpoint flattened rows with the wrong stride; heap overflow on any non-square grid \\
\svegp{18} & \code{independent()}/\code{dependent()} unguarded; null deref and heap overflow \\
\svegp{19} & \code{forward\_elimination()}/\code{reverse\_elimination()} never eliminated anything \\
\svegp{20} & uninitialised \code{Vertex*} handed to \code{add\_edge()} \\
\svegp{21} & \code{Process::mem\_size} never initialised \\
\svegp{22} & the partition count was not observable from the public API \\
\svegp{23} & the elimination direction was not selectable \\
\svegp{24} & the memory accounting never counted adjacency; every edge under-charged by more than half \\
\svegp{25} & one example's budget was a multiplier of a stale, undocumented unit \\
\svegp{26} & two parallel global states that had to agree; replaced by one tape object \\
\svegp{27} & \code{Process} had no destructor; a tape released other than through \code{finalize()} leaked its graph \\
\svegp{28} & the RAII surface: \code{Tape} and \code{Jacobian} \\
\svegp{29} & a tape re-entered from inside its own section restored the wrong predecessor \\
\svegp{30} & the adjacency arena \\
\bottomrule
\end{longtable}
\end{center}

\noindent
Still open by choice: \svegp{10}, an API call-order state machine that was
written and never wired in. An out-of-order call is diagnosed but not made
unrepresentable, except where the \code{Tape} surface makes it so.

\chapter{Measurement summary}

All figures GCC 13.3.0, \code{-O2}, x86-64.

\section{Graph representation}

\begin{center}
\begin{tabular}{@{}lrrrr@{}}
\toprule
 & original & edge arena & sorted adjacency & adjacency arena \\
\midrule
allocations / tape    & 39\,462 & 27\,908 & 17\,941 & $-72\%$ further \\
peak heap (bytes)     & 1\,780\,720 & 1\,559\,216 & 820\,896 & $-7$ to $-15\%$ further \\
\code{sizeof(Vertex)} & 120 & 120 & 72 & 72 \\
accounting ratio      & 2.44$\times$ & 2.14$\times$ & 1.63$\times$ & 1.51--1.56$\times$ \\
\bottomrule
\end{tabular}
\end{center}

\section{Cost of chunking}

16 variables, 120 steps, $\approx$1\,920 recorded operations.

\begin{center}
\begin{tabular}{@{}lrrr@{}}
\toprule
budget & partitions & time & replay share \\
\midrule
no AD at all & --- & 0.0085\,ms & --- \\
whole tape   & 1   & 5.78\,ms   & --- \\
200\,kB      & 8   & 15.40\,ms  & 62\% \\
50\,kB       & 29  & 18.35\,ms  & 68\% \\
12\,kB       & 121 & 36.40\,ms  & 84\% \\
\bottomrule
\end{tabular}
\end{center}

\section{Parallel independent tapes, 2 cores}

\begin{center}
\begin{tabular}{@{}lrrr@{}}
\toprule
budget & serial & parallel & speedup \\
\midrule
whole tape & 90.6\,ms & 50.7\,ms & 1.79$\times$ \\
200\,kB & 247.3\,ms & 129.0\,ms & 1.92$\times$ \\
50\,kB & 308.4\,ms & 157.4\,ms & 1.96$\times$ \\
12\,kB & 561.7\,ms & 281.1\,ms & 2.00$\times$ \\
\bottomrule
\end{tabular}
\end{center}

\section{Suite status}

\begin{center}
\begin{tabular}{@{}ll@{}}
\toprule
example & result \\
\midrule
\code{invariants} & PASS --- 12 structural values pinned; drift 1.51--1.56$\times$ \\
\code{fdcheck}    & 0 failures of 68 \\
\code{regress}    & PASS, 0 failures \\
\code{bratu}      & 13 chunks \\
\code{heat}       & 107 chunks \\
\code{heat\_assim}& PASS --- rms $1.001\times10^{-13}$ over 19 controls \\
\code{hybrdj1}    & 18 chunks \\
\code{pinn}       & PASS --- 2\,502 tapes \\
\code{test\_xor}  & PASS --- 4\,001 tapes, 8 chunks each \\
\bottomrule
\end{tabular}
\end{center}

\noindent
Zero warnings from a clean build under \code{-Wall -Wextra}; clean under
AddressSanitizer and LeakSanitizer.

\chapter{Known limitations}

\paragraph{No second derivatives.} A Hessian requires a tape whose edge weights
are themselves \code{active}, which requires the scalar type to be templated.
This is the largest change remaining and is not begun. Where an exact Hessian is
needed today, difference the AD gradient --- \code{examples/heat\_assim} does
exactly that, and its symmetry check agrees to $1.4\times10^{-11}$.

\paragraph{No parallelism within a tape.} Chapter~\ref{ch:threads} gives the
three reasons. Independent tapes parallelise; partitions do not.

\paragraph{The copy cost of \code{active} is undocumented in the header.} It is
described in this report and mitigated for independents, but a user who writes
the obvious growing-container pattern still pays silently.

\paragraph{Two counters share a name.} \code{Process::run\_counter} (productive
re-runs) and \code{TapeState::run\_counter} (replay position) are different
quantities with the same name inside what is now one object. Not a defect; a
consolidation waiting to happen.

\paragraph{\code{APICallState} is still not wired in.} An out-of-order
free-function call is diagnosed at runtime rather than made impossible. The
\code{Tape} surface makes most such orderings unrepresentable, which is the
better answer, but the free functions remain public.

\end{document}
